\subsection{CARTAN SUBALGEBRAS OF SEMI-SIMPLE LIE ALGEBRAS}

THEOREM 2. Assume that k is of characteristic 0. Let g be a semi-simple Lie algebra, h a Cartan subalgebra of g. Then h is commutative, and all of its elements are semi-simple in g (Chap. I, §6, no. 3, Def. 3).

Since \(\mathfrak{h} = \mathfrak{g}^0 (\mathfrak{h})\) , \(\mathfrak{h}\) is reductive (§1, Prop. 11), hence commutative since it is nilpotent. On the other hand, the restriction of the adjoint representation of \(\mathfrak{g}\) to \(\mathfrak{h}\) is semi-simple (loc. cit.), so the elements of \(\mathfrak{h}\) are semi-simple in \(\mathfrak{g}\) (Algebra, Chap. VIII, §5, no. 1).

COROLLARY 1. If \(x \in \mathfrak{h}\) and \(y \in \mathfrak{g}^{\lambda}(\mathfrak{h})\) , we have \([x, y] = \lambda(x)y\) .

Indeed, \(\mathfrak{g}^{\lambda (x)}(x) = \mathfrak{g}_{\lambda (x)}(x)\) since ad \(x\) is semi-simple.

COROLLARY 2. Every regular element of g is semi-simple.

Indeed, such an element belongs to a Cartan subalgebra (no. 3, Th. 1 (i)).

COROLLARY 3. Let \(\mathfrak{h}\) be a Cartan subalgebra of a reductive Lie algebra \(\mathfrak{g}\) . a) \(\mathfrak{h}\) is commutative.

\begin{enumerate}
\def\labelenumi{\alph{enumi})}
\setcounter{enumi}{1}
\tightlist
\item
  If \(\rho\) is a finite dimensional semi-simple representation of \(\mathfrak{g}\) , the elements of \(\rho(\mathfrak{h})\) are semi-simple.
\end{enumerate}

Let \(\mathfrak{c}\) be the centre of \(\mathfrak{g}\) , and \(\mathfrak{s}\) its derived algebra. Then \(\mathfrak{g} = \mathfrak{c} \times \mathfrak{s}\) , so \(\mathfrak{h} = \mathfrak{c} \times \mathfrak{h}'\) , where \(\mathfrak{h}'\) is a Cartan subalgebra of \(\mathfrak{s}\) (Prop. 2). In view of Th. 2, \(\mathfrak{h}'\) is commutative, hence so is \(\mathfrak{h}\) . Moreover, \(\rho(\mathfrak{h}')\) consists of semi-simple elements and so does \(\rho(\mathfrak{c})\) (Chap. I, §6, no. 5, Th. 4); assertion \(b\) ) follows.

